\subsubsection{Nivel 1: compra bajo incertidumbre}
\label{sec:modelo-compra}

El problema se descompone en dos niveles: la cantidad de plantas de cada especie que se adquiere en cada polígono y su acomodo en la retícula de plantación. Aquí se formula el primero, cuyo resultado es el dato de entrada del segundo (\cref{sec:modelo-espacial}). La compra se decide sin saber cuántas plantas nuevas sobrevivirán, por lo que se plantea como un programa estocástico con restricciones de probabilidad.

\paragraph{Datos y parámetros}
\label{subsec:compra-parametros}

Sea $\mathcal{S}$ el conjunto de las diez especies e $i$ un polígono de superficie $A_i$ (ha). Para cada $s\in\mathcal{S}$, $n_s^0$ es la densidad del plan de reforestación ($\sum_sn_s^0=658$ plantas/ha), $c_s$ el costo unitario de la planta sin logística y $q_s$ la supervivencia del primer año, observada en el monitoreo. $E_{is}$ es el número de plantas existentes de la especie $s$ en el polígono $i$ según los conteos de INECOL; en conjunto equivalen a 128.8 plantas/ha, 19.6\,\% de la meta, sobre los 183.01~ha de los 30 polígonos. La \cref{tab:compra-parametros} reporta estos valores. Se consideran además cuatro parámetros de política: el presupuesto por hectárea $B$, fijado por el socio formador en \$12{,}500; el sobrecosto logístico $\theta\in[0.05,0.10]$, como fracción del costo de la planta; la desviación máxima por especie respecto al plan, $\delta=0.10$; y el riesgo admisible de incumplimiento, $\alpha=0.05$.

\begin{table*}[t]
\centering
\caption{Parámetros por especie. $\bar{E}_s=\sum_iE_{is}/\sum_iA_i$: densidad media de plantas existentes; $c_s/q_s$: costo por planta sobreviviente.}
\label{tab:compra-parametros}
\small
\begin{tabular}{@{}lrrrrr@{}}
\toprule
Especie & $n_s^0$ (plantas/ha) & $c_s$ (\$) & $q_s$ & $\bar{E}_s$ (plantas/ha) & $c_s/q_s$ (\$) \\
\midrule
\textit{Agave lechuguilla}        &  42 & 26.0 & 0.943 &  8.1 & 27.6 \\
\textit{Agave salmiana}           & 196 & 26.0 & 0.959 & 38.0 & 27.1 \\
\textit{Agave scabra}             &  42 & 26.0 & 0.825 &  8.0 & 31.5 \\
\textit{Agave striata}            &  42 & 26.0 & 0.952 &  7.7 & 27.3 \\
\textit{Opuntia cantabrigiensis}  &  49 & 18.0 & 0.927 & 10.0 & 19.4 \\
\textit{Opuntia engelmannii}      &  38 & 21.0 & 0.899 &  8.0 & 23.4 \\
\textit{Opuntia robusta}          &  73 & 18.0 & 0.958 & 15.0 & 18.8 \\
\textit{Opuntia streptacantha}    &  64 & 18.0 & 0.886 & 12.5 & 20.3 \\
\textit{Prosopis laevigata}       &  86 & 26.5 & 0.821 & 16.5 & 32.3 \\
\textit{Yucca filifera}           &  26 & 26.0 & 0.884 &  5.0 & 29.4 \\
\midrule
Total                             & 658 &      &       & 128.8 &     \\
\bottomrule
\end{tabular}
\end{table*}

\paragraph{Formulación}
\label{subsec:compra-formulacion}

Sean $n_{is}\in\mathbb{Z}_{\geq0}$ la meta de plantas de la especie $s$ en el polígono $i$ y $x_{is}\in\mathbb{Z}_{\geq0}$ la cantidad a adquirir. La única variable aleatoria es el número de sobrevivientes de lo plantado, $S_{is}(x)\sim\operatorname{Bin}(x,q_s)$ (\cref{eq:binomial}), bajo independencia entre individuos. La \cref{tab:compra-notacion} resume la notación. Para cada polígono se resuelve
\begin{align}
\min_{n,\,x}\quad & (1+\theta)\sum_{s}c_s\Bigl(x_{is}+\mathbb{E}\!\left[D_{is}\right]\Bigr) \label{eq:obj}\\
\text{s.a.}\quad & \sum_{s}n_{is}=N_i, \label{eq:total}\\
& \left|n_{is}-n_s^0A_i\right|\leq\delta\,n_s^0A_i, && s\in\mathcal{S}, \label{eq:banda}\\
& (1+\theta)\sum_{s}c_sx_{is}\leq B\,A_i, \label{eq:presupuesto}\\
& \mathbb{P}\!\left(E_{is}+S_{is}(x_{is})\geq n_{is}\right)\geq1-\alpha, && s\in\mathcal{S}, \label{eq:probabilidad}
\end{align}
donde $N_i$ es $658A_i$ redondeado al entero más cercano y $D_{is}=\left(n_{is}-E_{is}-S_{is}(x_{is})\right)^{+}$ es el déficit al término del primer año. El objetivo \eqref{eq:obj} suma el costo de adquisición y el costo esperado de reponer el déficit, que representa la obligación contractual de reponer la planta muerta; la reposición se valora al mismo costo que la compra, logística incluida. La restricción \eqref{eq:total} fija la densidad total; \eqref{eq:banda} acota la composición por especie; \eqref{eq:presupuesto} limita el gasto, y \eqref{eq:probabilidad} exige que la meta de cada especie se alcance con probabilidad de al menos $1-\alpha$ \autocite{charnes1959}.

\begin{table*}[t]
\centering
\caption{Notación del modelo de compra.}
\label{tab:compra-notacion}
\small
\begin{tabularx}{\textwidth}{@{}l>{\RaggedRight\arraybackslash}p{0.17\textwidth}>{\RaggedRight\arraybackslash}X>{\RaggedRight\arraybackslash}p{0.2\textwidth}@{}}
\toprule
Símbolo & Tipo & Descripción & Unidad o valor \\
\midrule
$\mathcal{S}$, $s$ & Conjunto e índice & Especies a plantar & 10 especies \\
$i$ & Índice & Polígono de reforestación & 30 polígonos \\
$A_i$ & Dato & Superficie del polígono $i$ & ha \\
$N_i$ & Dato & Meta total del polígono, $658A_i$ redondeado & plantas \\
$n_s^0$ & Dato & Plantas de la especie $s$ por hectárea según el plan & plantas/ha \\
$c_s$ & Dato & Costo unitario de la planta, sin logística & \$/planta \\
$E_{is}$ & Dato & Plantas existentes de la especie $s$ en el polígono $i$ & plantas \\
$q_s$ & Parámetro estimado & Probabilidad de que una planta nueva sobreviva el primer año & $[0,1]$ \\
$B$ & Parámetro de política & Presupuesto por hectárea & \$12{,}500 \\
$\theta$ & Parámetro de política & Sobrecosto logístico como fracción del costo de la planta & $[0.05,\,0.10]$ \\
$\delta$ & Parámetro de política & Desviación relativa máxima de la meta de cada especie & 0.10 \\
$\alpha$ & Parámetro de política & Probabilidad admisible de no alcanzar la meta de una especie & 0.05 \\
$S_{is}(x)$ & Variable aleatoria & Sobrevivientes al primer año de $x$ plantas nuevas; $\operatorname{Bin}(x,q_s)$ & plantas \\
$D_{is}$ & Variable aleatoria & Déficit al término del primer año, $(n_{is}-E_{is}-S_{is})^{+}$ & plantas \\
$n_{is}$ & Variable de decisión & Meta de plantas de la especie $s$ en el polígono $i$ & plantas; entera \\
$x_{is}$ & Variable de decisión & Plantas de la especie $s$ a adquirir y plantar en el polígono $i$ & plantas; entera \\
\bottomrule
\end{tabularx}
\end{table*}

Antes de resolver conviene estimar el orden de magnitud. Con la composición del plan, el costo de cubrir el faltante, $\sum_sc_s(n_s^0-\bar{E}_s)$, asciende a \$12{,}456/ha, y a \$13{,}683/ha al corregirlo por la mortalidad del primer año, $\sum_sc_s(n_s^0-\bar{E}_s)/q_s$, aun sin margen de seguridad ni logística. Con $B=\$12{,}500$, la restricción \eqref{eq:presupuesto} es por tanto prácticamente activa, y la factibilidad depende de la banda \eqref{eq:banda}, que permite desplazar la composición hacia especies con menor costo por planta sobreviviente (\$18.8 para \textit{O.~robusta} frente a \$32.3 para \textit{P.~laevigata}). Cuando el problema es infactible, el modelo entrega el presupuesto mínimo $B_i^{*}$ con el que dejaría de serlo.

\paragraph{Supuestos y extensiones}
\label{subsec:compra-supuestos}

El modelo descansa en cuatro supuestos. (i) La supervivencia es independiente entre individuos; los eventos que afectan a toda la plantación, como la sequía, aumentan la varianza de $S_{is}$ y pueden incorporarse sorteando por escenario una tasa común $q_s^{(k)}$, aunque un solo año de monitoreo no permite estimar su dispersión. (ii) Los conteos de plantas existentes son exactos y esas plantas no mueren en el horizonte. Si el conteo fuera incierto, $E_{is}$ puede sortearse en cada escenario como $\operatorname{Poisson}(\bar{E}_sA_i)$: el estadístico de dispersión de Pearson de los conteos respecto a $\bar{E}_sA_i$ está entre 0.58 y 1.45 para todas las especies, sin evidencia de la sobredispersión que justificaría una binomial negativa \autocite{cameron2013}. (iii) El horizonte es de un año, aunque el monitoreo muestra que la mortalidad continúa el segundo; para obligaciones de reposición multianuales, el término de reposición se extiende con la cadena de Markov por posición de la \cref{subsec:markov}. (iv) El costo logístico es proporcional al costo de la planta; los costos fijos por viaje o por vivero se incorporan con variables binarias de activación sin alterar la estructura del modelo.
